\section{Preliminaries: a fractional reference schedule}
\label{sec:foundations-lp}

We begin with the time-indexed relaxation underlying Li's common-quantile
approach \cite[Sections~2.1 and~3]{Li2017}. We spell out the relaxation and the
interpolation used in this paper, since the latter also fixes the units and
boundary conventions used in the proof.

Let $\mathcal O=\{1,\ldots,n\}$ be the set of original jobs.
Write $i\to j$ when the input graph contains an arc from $i$ to $j$, and
call $i$ a predecessor of $j$ when $i\ne j$ and there is a directed path
from $i$ to $j$.
For a finite set of jobs, vertical bars denote its cardinality. For a
measurable subset of the line or plane, they denote its length or area,
respectively (Lebesgue measure). The \emph{depth}
$a_j$ is the largest number of jobs on a directed path ending at $j$,
including $j$. A set of jobs is \emph{predecessor closed} if it contains
every predecessor of each of its members. A \emph{chain} is a sequence of
jobs ordered by the precedence relation, and the \emph{height} of a set
of jobs is the largest number of its jobs in a chain (zero for the empty
set). Every schedule satisfies $C_j\ge a_j$: the jobs on such a
path must execute one after another. The depths can be computed in polynomial
time by repeatedly processing a job whose predecessors have been processed.
For an integer $s\ge1$, the \emph{slot} $s$ is
the interval $(s-1,s]$. If $m\ge n$, schedule each job $j$ in
slot $a_j$. This is feasible and achieves all these lower bounds. Thus we
may assume $1\le m<n$ below.

\begin{lemma}[A finite integer horizon]
There is an optimal schedule in which every job completes at an integer
time between $1$ and $n$.
\end{lemma}
\begin{proof}
Consider any feasible schedule and order the jobs assigned to each machine
by their execution times. Add an arc between consecutive jobs in each
machine order. Along each added or original arc, the second job starts no
earlier than the first job's completion. Completion times therefore increase
strictly along arcs, so the augmented graph is acyclic. Let $a'_j$ be the number
of vertices on a longest augmented path ending at $j$. Execute $j$ in slot
$a'_j$. The added arcs ensure that jobs
assigned to the same machine do not overlap, and the original arcs ensure
that precedence is respected. Every path has at most $n$ jobs, so
$a'_j\in\{1,\ldots,n\}$. Along each augmented path ending at $j$, unit
processing and the order constraints imply that the original completion
time of $j$ is at least the length of the path. Therefore the new completion
time $a'_j$ is no larger than the original one. Since weights are
nonnegative, the transformed schedule costs no more. There are only finitely
many integer-slot schedules within this horizon, so one of minimum cost
among them is optimal even when compared with arbitrary feasible schedules.
\end{proof}

For each job $j$ and slot $s\in\{1,\ldots,n\}$, the variable
$x_{j,s}$ is its fractional completion mass in slot $s$. We use the
following linear program (LP). A sum with an upper limit below its lower
limit is zero.
\begin{align*}
 \min\quad &\sum_{j\in\mathcal O}w_j\cdot\sum_{s=1}^n s\cdot x_{j,s},\\
 \sum_{s=1}^n x_{j,s}&=1 &&\text{for every }j,\\
 \sum_{j\in\mathcal O} x_{j,s}&\le m &&\text{for every }s,\\
 \sum_{s'=1}^{s}x_{j,s'}&\le\sum_{s'=1}^{s-1}x_{i,s'}
       &&\text{for every }i\to j\text{ and every }s,\\
 x_{j,s}&\ge0.
\end{align*}
The precedence constraint says that mass completing $j$ by slot $s$ must
be matched by mass completing each job with an arc to $j$ by slot $s-1$. For a
feasible solution, put $C_j^{\mathrm{LP}}=\sum_{s=1}^n s\cdot x_{j,s}$.

The depth-support observation follows the corresponding argument in
Li's analysis \cite[Claim~3.1]{Li2017}.

\begin{lemma}[Lower bound and depth support]
\label{lem:lp-support}
Every feasible integer-slot schedule completing by time $n$ gives a feasible
LP solution with the same objective. Moreover, every feasible LP solution satisfies
$x_{j,s}=0$ for $s<a_j$.
\end{lemma}
\begin{proof}
Put unit mass in the completion slot of each job. Assignment, capacity,
and precedence are then immediate, and the objective agrees with the
schedule. Hence the LP optimum is a lower bound on the scheduling optimum.
For the support claim, take a path $i_1\to\cdots\to i_{a_j}=j$. Repeatedly
applying the cumulative precedence inequalities gives, for $s<a_j$,
\[
 \sum_{s'=1}^{s}x_{j,s'}
 \le \sum_{s'=1}^{s-(a_j-1)}x_{i_1,s'}=0.
\]
Nonnegativity now gives the claim. The same repeated application shows how the constraints on
arcs propagate along any path.
\end{proof}

The analysis below applies to any feasible LP solution. The algorithm uses
an optimal solution, so its fractional objective is a lower bound on the
optimal scheduling cost.

\subsection{Virtual capacity padding}
\label{sec:foundations-padding}

It is useful to make the LP capacity exactly full in the analysis.
Summing the assignment constraints gives
\[
 \sum_{s=1}^n\left(m-\sum_{j\in\mathcal O}x_{j,s}\right)=n\cdot(m-1).
\]
If $m>1$, introduce a set $\mathcal D$ of $n\cdot(m-1)$ independent
\emph{virtual jobs}, each of weight zero and with no precedence arcs.
Give each virtual job the fractional distribution
\[
 x_{i,s}=\frac{m-\sum_{j\in\mathcal O}x_{j,s}}{n\cdot(m-1)}
 \qquad(i\in\mathcal D,\ 1\le s\le n).
\]
If $m=1$, let $\mathcal D$ be empty.

\begin{lemma}[Exact fractional capacity]
The virtual distributions are valid, and after padding every slot has
exactly $m$ units of completion mass. The original-job distributions and the
weighted LP objective are unchanged.
\end{lemma}
\begin{proof}
By the capacity constraints, every numerator is nonnegative. For $m>1$,
the displayed sum shows that each virtual job's masses sum to one and
their total mass in each slot fills the unused capacity. For $m=1$, the
nonnegative unused capacities sum to zero, so every slot was already full.
Virtual jobs have zero weight and do not alter any original distribution.
\end{proof}

Let $\mathcal J=\mathcal O\cup\mathcal D$. Thus
$|\mathcal J|=m\cdot n$ and
\[
 \sum_{i\in\mathcal J}x_{i,s}=m\qquad(1\le s\le n).
\]
The algorithm will schedule only the original jobs. The later
virtual-pool transfer proposition explains why padded quantities may
nevertheless be used to bound that schedule.

\subsection{Continuous completion quantiles}
\label{sec:foundations-interpolation}

Spread mass $x_{i,s}$ uniformly over the slot $(s-1,s]$. Its \emph{cumulative distribution function} (CDF) is, for $t\in[s-1,s]$,
\[
 F_i(t)=\sum_{s'=1}^{s-1}x_{i,s'}+(t-s+1)\cdot x_{i,s}.
\]
Extend $F_i$ by zero for $t\le0$ and by one for $t\ge n$. Each $F_i$ is
continuous and nondecreasing. Exact slot capacity gives the identity
\begin{equation}
 \sum_{i\in\mathcal J}F_i(t)=m\cdot t\qquad(0\le t\le n).
 \label{eq:foundations-continuous-capacity}
\end{equation}
Indeed, all earlier slots contribute $m$ each and the current slot
contributes $m\cdot(t-s+1)$.

\begin{lemma}[Continuous precedence]
For every arc $i\to j$ and every real $t$,
\begin{equation}
 F_j(t)\le F_i(t-1).
 \label{eq:foundations-continuous-precedence}
\end{equation}
\end{lemma}
\begin{proof}
At integer times inside the horizon, this is the cumulative LP precedence
constraint. At the remaining integer times it follows from the extensions
and the constraint at the horizon. Between successive integers, both sides
are affine, meaning linear plus a constant, so the difference is affine as
well. It is nonpositive at both
endpoints and therefore throughout the interval.
\end{proof}

For $0<\alpha<1$ define the \emph{generalized inverse}, or \emph{quantile},
\[
 M_i(\alpha):=\inf\{t\in[0,n]:F_i(t)\ge\alpha\}.
\]
At the endpoints, extend $M_i$ by its one-sided limits:
\[
 M_i(0)=\lim_{\alpha\downarrow0}M_i(\alpha),\qquad
 M_i(1)=\lim_{\alpha\uparrow1}M_i(\alpha).
\]
These limits exist because $M_i$ is nondecreasing and bounded. The
extension at $0$ is not obtained by substituting $\alpha=0$ into the
infimum above. Intuitively, $M_i(\alpha)$ is the interpolated time by
which an $\alpha$-fraction of job $i$ has completed. Although $F_i$ is
continuous, $M_i$ can jump when $F_i$ has a flat interval.

For a statement $A$, write $\mathbf1\{A\}$ for its indicator: one when $A$
holds and zero otherwise.

\begin{lemma}[Quantile separation and mean]
For every $0<\alpha<1$, continuity gives $F_i(M_i(\alpha))=\alpha$.
For each arc $i\to j$,
\begin{equation}
 M_i(\alpha)+1\le M_j(\alpha).
 \label{eq:foundations-quantile-separation}
\end{equation}
Also $M_j(0)\ge a_j-1$, and
\begin{equation}
 \int_0^1M_j(\alpha)\,d\alpha
   =\sum_{s=1}^n(s-\tfrac12)\cdot x_{j,s}
   =C_j^{\mathrm{LP}}-\tfrac12.
 \label{eq:foundations-quantile-mean}
\end{equation}
\end{lemma}
\begin{proof}
Continuity of $F_i$ and the definition of its generalized inverse give
$F_i(M_i(\alpha))=\alpha$. Applying \eqref{eq:foundations-continuous-precedence} at
$t=M_j(\alpha)$ gives $F_i(M_j(\alpha)-1)\ge\alpha$ and hence the
separation inequality. Depth support puts all of $j$'s mass in slots
$a_j$ and later, so $M_j(\alpha)\ge a_j-1$.

For the mean, interchange the order of integration in
$\int_0^1\int_0^n\mathbf1\{M_j(\alpha)>t\}\,dt\,d\alpha$. For fixed $t$,
the measure of $\alpha$ with $M_j(\alpha)>t$ is $1-F_j(t)$, giving
\[
 \int_0^1 M_j(\alpha)\,d\alpha
 =\int_0^n(1-F_j(t))\,dt.
\]
This is the mean of the interpolated distribution. The portion in slot $s$
has mass $x_{j,s}$ and mean $s-\tfrac12$, which proves the identity. The
half-unit correction will be kept throughout the analysis. The LP objective
itself still uses integer completion slots.
\end{proof}
